\documentclass[11pt]{article}
\usepackage[margin=1in]{geometry}
\usepackage{amsmath,amssymb,amsthm,booktabs,microtype,hyperref}
\hypersetup{hidelinks,pdfauthor={T},pdftitle={When Abstention Fails to Transfer Across an Ethical Benchmark Shift}}
\newtheorem{theorem}{Theorem}
\title{When Abstention Fails to Transfer Across an Ethical Benchmark Shift}
\author{T}
\date{}
\begin{document}
\maketitle
\begin{abstract}
An assistant may defer cases on which a learned judgment is uncertain.
That strategy is useful only if confidence continues to rank errors
when the cases change. We prove a sharp transfer bound for risk
conditional on a fixed acceptance rule. When target coverage is
observed, an exact interval can be tighter. We then fit a lexical classifier
on the public ETHICS commonsense corpus and freeze three abstention
rules on a development split. At the strictest raw-confidence threshold,
accuracy among retained cases rises on the ordinary test but falls on
the hard test. The rule also omits every short test case at that
threshold. Reporting coverage and error separately exposes both effects.
At equal short- and long-text quotas, three confidence rankings
beat a score-free baseline on the ordinary split. None beats it
on the hard split.
\end{abstract}

\section*{Summary and scope}
\paragraph{Question.} When can a frozen abstention rule make accepted
ethical judgments less reliable after a shift?
\paragraph{Finding.} A sharp total-variation bound applies at fixed
acceptance. A sharper interval uses the observed coverage. On ETHICS,
the strict raw-confidence rule retains zero short
texts, and reweighting retained texts by length and label does not
explain the hard split's accuracy drop.
\paragraph{Evidence and scope.} The lexical
experiment is separate empirical evidence on public dataset labels. The
result supports checking retained errors and omissions before automating
actions. It gives no estimate of language-model behavior or real-world
moral correctness.

\section{Abstention changes the population being judged}
A tool-use assistant may ask for human review when a classifier
cannot decide whether an action is permissible. When the classifier
accepts a case, the relevant risk is conditional on acceptance.
A lower acceptance rate can improve that risk when uncertainty
is informative. It can also concentrate mistakes or leave one
type of user request entirely without automatic service.

Selective classification formalizes the coverage and error
tradeoff \cite{geifman}. Empirical work on uncertainty under
distribution shift shows the need to test confidence on changed
inputs \cite{ovadia}. Here we use a deliberately small classifier
to isolate what a frozen abstention threshold does on two public
splits of an ethical judgment benchmark \cite{hendrycks}.

\section{A sharp conditional transfer bound}
Let $P$ and $Q$ be two laws on the joint input and label.
Fix a classifier and acceptance event $A$ before evaluating
either law. Let $E$ be the classifier's error event,
$p=P(A)>0$, $q=Q(A)>0$, and
$r_P=P(E\mid A)$, $r_Q=Q(E\mid A)$.

\begin{theorem}
If $\mathrm{TV}(P,Q)\le\varepsilon$, then
\[
 |r_Q-r_P|\le\min\!\left\{1,\frac{\varepsilon}{\max(p,q)}\right\}.
\]
\end{theorem}
\begin{proof}
Write $\Delta_1=Q(A\cap E)-P(A\cap E)$ and
$\Delta_0=Q(A\cap E^c)-P(A\cap E^c)$. Both have absolute
value at most $\varepsilon$. Direct expansion gives
\[
q(r_Q-r_P)=(1-r_P)\Delta_1-r_P\Delta_0 .
\]
The absolute value is at most $\varepsilon$, since the two
nonnegative coefficients sum to one. Reversing the roles of
$P,Q$ gives the same bound with denominator $p$. Combining
both and using the unit range proves the claim.
\end{proof}

For sharpness, use three states for retained error, retained
correct, and rejection. Let $P=(0,c,1-c)$ and
$Q=(\varepsilon,c-\varepsilon,1-c)$ with
$0<\varepsilon\le c\le1$. Both have coverage $c$ and their
conditional risks differ by $\varepsilon/c$. Exact enumeration
on denominator-eight laws checked all 1,936 ordered pairs with
positive coverage and found 560 nontrivial equalities.
The theorem concerns joint input and label laws. Text similarity
alone cannot certify $\varepsilon$. We do not estimate a
deployment shift radius from the corpus.

\section{Using target coverage before labels arrive}
Coverage can be measured before anyone adjudicates accepted
predictions. Let $a=P(A\cap E)$ and write $d=q-p$. Assume
that each of the three cells for accepted error, accepted correct,
and rejection can receive mass in the target model. Set
\begin{align*}
 L(q)&=\max\{0,a+[d]_+-\varepsilon\},\\
 U(q)&=\min\{q,a+\min(0,d)+\varepsilon\},
\end{align*}
where $[d]_+=\max(0,d)$.

\begin{theorem}
If $|q-p|>\varepsilon$, the coverage is incompatible with
the assumed total-variation bound. Otherwise the sharp range
of retained risk is $[L(q)/q,U(q)/q]$. If target coverage is known
only to lie in a positive closed interval, first intersect the
interval with $[p-\varepsilon,p+\varepsilon]$. If the intersection
is nonempty, the coverage nearest $p$ gives the sharp endpoints
over that entire intersection.
\end{theorem}
\begin{proof}
The changes in the three cell masses are $x-a$, $d-(x-a)$,
and $-d$, where $x=Q(A\cap E)$. They sum to zero. For three
such changes the total variation is their largest absolute
value, since the sum of the positive changes equals the
absolute value of the sum of the negative changes. Bounding
all three and imposing $0\le x\le q$ gives $[L(q),U(q)]$.
Every allowed triple of masses attains its projected total
variation by retaining the source law conditional on each cell
of positive source mass. For a source-null cell, use any law
supported there. This gives the sharp range in the stated model.

For $q\le p$, the lower ratio is
$\max(0,(a-\varepsilon)/q)$ and is nonincreasing in $q$.
The upper ratio is $\min(1,1+(a+\varepsilon-p)/q)$ and
is nondecreasing. For $q\ge p$, the lower ratio is
$\max(0,1+(a-\varepsilon-p)/q)$ and is nondecreasing,
and the upper ratio is $\min(1,(a+\varepsilon)/q)$ and is
nonincreasing. Thus the coverage nearest $p$ gives the widest
risk interval among the feasible coverages.
\end{proof}

For example, suppose source coverage is $0.60$ with retained
risk $0.20$. A shift radius of $0.20$ and target coverage of
$0.40$ give the exact retained-risk interval $[0,0.30]$.
The earlier symmetric inequality gives an upper bound of $8/15$.
The sharper result uses the direction of the coverage change.
The verifier agreed with direct enumeration on 3,168 fixed
coverage cases and 14,256 interval cases. A measured coverage
outside the feasible interval challenges the radius assumption.
Coverage inside the interval still leaves the radius unverified.

\section{Frozen comparison on public text}
The public ETHICS commonsense files contain short and long
scenarios with binary judgment labels. The authors' hard split
was produced by adversarial filtering \cite{hendrycks}. We
use their split as a stress condition. Short and long are
dataset strata, not demographic groups.

Normalize and hash every text. Drop conflicting duplicates and
keep one copy of each remaining normalized text. Assign training
texts to fitting or development by hash modulo five. This gives
11,093 fitting texts and 2,816 development texts. Remove any
normalized training digest from each official test split. This
removes 1,160 ordinary test texts, leaving 2,725. The hard split
has 3,964 remaining texts and no exact normalized training
overlap. This filter controls exact normalized overlap only.
The same public splits have been used for earlier exploratory
baselines, so this is a frozen comparison within an existing
benchmark rather than a new external blind evaluation.

Fit one multinomial word classifier with additive smoothing.
Its 14,269-word vocabulary is selected on the fitting texts.
Let $s$ be its log odds and predict the positive label when
$s>0$. For abstention, compare $|s|$, $|s|$ divided by the
number of known tokens, and separate $|s|$ thresholds for
short and long texts. The prediction itself is identical
across the three rules. Calibrate cutoffs on development data
to retain 80, 50, or 20 percent, keeping ties. Apply every
threshold unchanged to both test splits. The release of these
thresholds is a coverage calibration, with no guarantee that
it controls error on a changed law.

\begin{center}
\small
\begin{tabular}{lrrrr}
\toprule
Rule and development target & \multicolumn{2}{c}{Ordinary} & \multicolumn{2}{c}{Hard} \\
 & Coverage & Retained accuracy & Coverage & Retained accuracy \\
\midrule
Forced prediction & 100.00 & 58.79 & 100.00 & 45.59 \\
Raw confidence, 50\% & 26.28 & 72.21 & 53.23 & 40.76 \\
Per-token confidence, 50\% & 62.83 & 59.99 & 47.05 & 45.63 \\
Group thresholds, 50\% & 48.99 & 63.00 & 47.33 & 43.28 \\
Raw confidence, 20\% & 8.40 & 89.52 & 18.84 & 37.08 \\
Per-token confidence, 20\% & 30.35 & 57.19 & 18.26 & 50.55 \\
Group thresholds, 20\% & 19.38 & 64.39 & 17.94 & 46.13 \\
\bottomrule
\end{tabular}

\medskip
\footnotesize Entries are percentages. The receipt includes the 80 percent target.
\end{center}
At the strictest raw threshold the ordinary split retained 229
texts, of which 205 were correct. The hard split retained 747
texts, of which 277 were correct. Both retained zero short
texts. Group-specific calibration retained both strata at this
target, but still achieved only 328 correct predictions among
711 retained hard texts. Normalizing the score changed the
coverage mix and gave 366 of 724 correct at its strictest
hard threshold. These rules trade off different coverage and
accuracy profiles, and the experiment does not select a winning
rule using the test results.

Consider errors in the direction of declaring a text permissible
when the dataset marks it as wrong. Among 429 hard texts declared
permissible by the strict raw threshold, the dataset labels
275 as wrong. Under forced prediction the corresponding counts are
1,265 of 2,246. The stricter rule makes fewer such errors in
absolute terms because it acts less often. The fraction among
its permissive predictions is higher. Both quantities matter
when a deployment uses the output to decide whether to act.

\section{Does the mix of retained texts explain the drop?}
Fix a rule and form cells by crossing the dataset's short and long
categories with its binary label. In split $D\in\{O,H\}$, let $w_{Dg}$
be the share of retained texts in cell $g$ and $a_{Dg}$ their accuracy.
On common support, the difference in retained accuracy is
\[
 \sum_g (w_{Hg}-w_{Og})\frac{a_{Hg}+a_{Og}}{2}
 +\sum_g \frac{w_{Hg}+w_{Og}}{2}(a_{Hg}-a_{Og}).
\]
The two terms separate changes in shares and changes in accuracy
within cells. This is the symmetric Kitagawa decomposition
\cite{oaxaca}. Expanding each summand cancels both cross terms and
leaves $w_{Hg}a_{Hg}-w_{Og}a_{Og}$.

A second comparison holds ordinary cell accuracies fixed and gives
them the hard split's shares, yielding $\sum_g w_{Hg}a_{Og}$.
If a hard cell has no ordinary observations, its term is bounded
between zero and its hard share. Every reweighting on ordinary
support must lie between the smallest and largest ordinary cell
accuracies, because the weights are nonnegative and sum to one.

\begin{center}
\small
\begin{tabular}{lrr}
\toprule
Retained cell at the strict raw threshold & Ordinary correct/total & Hard correct/total \\
\midrule
Long, label 0 & $118/132$ & $154/349$ \\
Long, label 1 & $87/97$ & $123/398$ \\
\bottomrule
\end{tabular}
\end{center}

Reweighting the ordinary accuracies to match the hard split yields
$89.55\%$, compared with the observed $37.08\%$ on the hard split.
Even assigning all weight to the lower-accuracy ordinary cell gives
$118/132$, or $89.39\%$. The symmetric decomposition assigns
$-0.71$ percentage points to changed shares and $-51.73$ points
to within-cell changes. The strict raw threshold is the only one
of the ten evaluated regimes whose hard accuracy falls below every
ordinary length-by-label cell. Across the ten regimes, applying the
hard shares to the ordinary accuracies gives a higher estimate
than the observed hard accuracy. The comparison is descriptive:
other features of the texts can change within a length-by-label
cell, and the calculation does not identify which feature drove
the errors.

\section{Comparing rankings at the same coverage}
Different cutoffs accept different numbers of short and long
texts. To test the ranking itself, we kept the same fitted
classifier and gave each rule the same acceptance quota within
each length group. A score-free rule orders normalized text
digests and serves as a baseline. We compared it with raw
confidence, confidence per known token, and a score learned
from development-set correctness.

Split the held development fifth into 1,439 cases for learning
the ranking and 1,377 for calibrating the cutoffs. Within each
length and predicted-label group, divide per-token confidence
into five bins using the ranking cases. Rank by the smoothed
correctness rate of the corresponding bin. The rate is the
number of correct predictions plus one, divided by the bin count
plus two. Break all ties by the same text digest order. These
rules were fixed before the new scoring run on the previously
explored public splits.

In the matched comparison, accept the top fifth of each length
group in each test batch, rounding up. This procedure uses the
unlabeled batch to set its quota. Flipping every test label
left every selected set unchanged. Each rule retained 422 short
and 124 long texts on the ordinary split and 341 short and
452 long texts on the hard split.
\begin{center}
\small
\begin{tabular}{lrr}
\toprule
Ranking & Ordinary correct of 546 & Hard correct of 793 \\
\midrule
Score-free digest & 320 & 348 \\
Raw confidence & 351 & 346 \\
Per-token confidence & 357 & 342 \\
Development correctness & 341 & 345 \\
\bottomrule
\end{tabular}
\end{center}

The per-token ranking gained 37 correct decisions over the
score-free rule on the ordinary split and lost six on the hard
split. Learning correctness on development data did not reverse
that shift. In the separate threshold-transfer comparison, the
cutoffs were fitted only on the calibration half and applied
unchanged. Every rule retained both length groups. Across the
four rules, hard-split acceptance counts ranged from 320 to 360
for short texts and from 379 to 487 for long texts. A batch quota
controls representation directly, while a fixed threshold offers
a streaming rule whose coverage can drift. Neither comparison
uses a new external holdout.

\section{Application and limits}
A service can use a confidence score to prioritize human review,
provided it checks both coverage and errors under the relevant
conditions. Keep authorization and high-impact constraints
outside the score. To test a proposed deferral rule, preregister
its threshold and evaluate both retained accuracy and omissions
in each affected group. A better follow-up would use independent
adjudication of actual agent actions on a new cohort and would
measure what happens to deferred cases. This study measures
agreement with dataset judgments for one small classifier.
It supplies no estimate of moral correctness, language-model
behavior, or the consequences of real tool use.

The reweighting and decomposition use exact rational arithmetic.

\paragraph{Review status.} Independent manuscript review is pending.

\begin{thebibliography}{9}
\bibitem{geifman} Yonatan Geifman and Ran El-Yaniv.
Selective Classification for Deep Neural Networks. NeurIPS, 2017.
\url{https://arxiv.org/abs/1705.08500}.
\bibitem{ovadia} Yaniv Ovadia et al. Can You Trust Your Model's
Uncertainty? Evaluating Predictive Uncertainty Under Dataset Shift.
NeurIPS, 2019. \url{https://arxiv.org/abs/1906.02530}.
\bibitem{hendrycks} Dan Hendrycks et al.
Aligning AI With Shared Human Values. ICLR, 2021.
\url{https://arxiv.org/abs/2008.02275}.
\bibitem{oaxaca} Ronald L. Oaxaca and Eva Sierminska.
Oaxaca-Blinder Meets Kitagawa: What Is the Link?
\emph{PLOS ONE} 20(5), e0321874, 2025.
\url{https://journals.plos.org/plosone/article?id=10.1371/journal.pone.0321874}.
\end{thebibliography}
\end{document}
